%% IEEE conference paper (ITSC / IV submission)
%% Compile on Overleaf with pdflatex (IEEEtran class is bundled).
\documentclass[conference]{IEEEtran}
\usepackage[utf8]{inputenc}
\usepackage[T1]{fontenc}
\usepackage{graphicx}
\usepackage{cite}
\usepackage{amsmath,amssymb,amsfonts}
\usepackage{algorithmic}
\usepackage{booktabs}
\usepackage{url}

\def\BibTeX{{\rm B\kern-.05em{\sc i\kern-.025em b}\kern-.08em T\kern-.1667em\lower.7ex\hbox{E}\kern-.125emX}}

\begin{document}

\title{Low-Latency Edge LLM Tactical Decision-Making for High-Speed Autonomous Overtaking}

\author{\IEEEauthorblockN{Author Name}
\IEEEauthorblockA{\textit{Department} \\
\textit{University}\\
City, Country \\
email@example.com}}

\maketitle

\begin{abstract}
Large language models (LLMs) have recently been used as decision-makers for
autonomous driving, but existing approaches rely on cloud APIs whose
multi-second inference latency precludes real-time, safety-critical use. We
argue that the LLM should act only as a \emph{low-frequency tactical} layer ---
issuing coarse overtake/hold intents --- while a classical 20~Hz controller
handles trajectory tracking and safety. Under this decomposition we fine-tune a
Qwen2.5-1.5B model with LoRA and show that a short-output variant reaches a
decision latency of \textbf{0.59~s} on a consumer 8~GB GPU, roughly an order of
magnitude faster than cloud-API baselines, while retaining 86.1\% decision
accuracy --- the Pareto-optimal point among 0.5B--7B alternatives. In
closed-loop CARLA experiments the system completes overtakes with 95--100\%
pass rate, 0--2\% collision, and median lateral clearance $\geq$2.2~m across ten
corner--lead scenarios, generalizing \emph{zero-shot} from training on only two
corner geometries. We additionally report two honest negative results:
fine-tuning on multi-scenario rule-expert data \emph{regresses} closed-loop
performance via catastrophic forgetting, and a CBF-QP safety filter
underperforms a simple geometric shield. Our results establish a quantified
latency--safety boundary for embedding LLMs in high-speed driving.
\end{abstract}

\begin{IEEEkeywords}
autonomous driving, large language models, low-latency inference, overtaking, hierarchical control, edge AI
\end{IEEEkeywords}

\section{Introduction}
\IEEEPARstart{A}{utonomous} driving has conventionally been organized as a
perception--planning--control stack, in which a motion planner produces a
trajectory that a low-level controller tracks at high frequency. This
decomposition is well understood and robust, but it struggles with situations
that demand semantic or commonsense reasoning: interpreting an erratic
neighbor's intent, deciding whether a narrowing gap is safe to enter, or
recognizing that a vehicle is deliberately blocking an overtake. Large language
models (LLMs) promise to close exactly this gap, because they encode broad
world knowledge and can articulate \emph{why} a maneuver is appropriate
\cite{sha2023languagempc,xu2023drivegpt4}. A growing line of work therefore
positions the LLM as a decision-maker and translates its reasoning into driving
actions \cite{sha2023languagempc}.

Yet this line shares a common bottleneck: the decision is produced by a cloud
model such as GPT-4, whose end-to-end latency --- autoregressive decoding over
hundreds of tokens plus a network round-trip --- is measured in seconds. This is
acceptable for offline annotation or slow scenarios, but not for high-speed
maneuvering. At 30~m/s a vehicle covers 30~m per second; a decision that arrives
three seconds late describes a traffic scene that no longer exists, and any
safety-critical intent (e.g., ``hold, do not enter the gap'') arrives after it is
already too late. Real-time deployment therefore requires the decision model to
be \emph{simultaneously accurate and fast}, a trade-off that has been asserted
but, to our knowledge, not quantified in a closed driving loop.

We make a simple architectural observation that relaxes the latency
requirement. Following the classical hierarchical decomposition, we confine the
LLM to a \emph{tactical layer} that runs at roughly 1~Hz and emits a coarse,
structured intent --- \emph{overtake or hold}, plus an attack offset and a speed
factor --- while a classical \emph{control layer} runs at 20~Hz and owns
trajectory tracking and safety. Under this decomposition, decision latency is
absorbed into the decision \emph{cadence} rather than the control loop: a
sub-second decision is tolerable, whereas a multi-second one is stale (we
quantify this boundary in Sec.~\ref{sec:stale}). This insight lets us replace
the cloud model with a \emph{small, locally fine-tuned} LLM, turning an
order-of-magnitude latency reduction from a nice-to-have into the enabling
condition for real-time use.

Concretely, we fine-tune Qwen2.5-1.5B with LoRA on a single consumer 8~GB GPU
and study high-speed overtaking on a closed CARLA ring. Our contributions are:
\begin{enumerate}
\item \textbf{A latency--accuracy Pareto characterization} across 0.5B--7B
models, showing that a short-output 1.5B variant is the optimal operating point
(0.59~s, 86.1\%), outperforming 3B (2.02~s, 71.3\%) and 7B (1.41~s, 69.6\%) on
\emph{both} axes, because the larger models cannot be adequately fine-tuned
within the memory budget.
\item \textbf{A closed-loop evaluation} in CARLA in which the system completes
overtakes with 95--100\% pass rate and median clearance $\geq$2.2~m across ten
corner--lead scenarios, generalizing \emph{zero-shot} from training on only two
corner geometries.
\item \textbf{A quantified stale-latency boundary} showing that 0.5~s of
decision delay is tolerated while 0.8~s degrades safety --- the quantitative
justification for the latency reduction.
\item \textbf{Two control-layer corrections and two honest negative results}
that are independently useful: an attack-geometry correction, an inside-line
hold, the failure of multi-scenario behavior cloning (catastrophic forgetting),
and the underperformance of a CBF-QP filter relative to a simple geometric
shield.
\end{enumerate}

\section{Related Work}
\subsection{LLMs for Autonomous Driving}
LanguageMPC \cite{sha2023languagempc} employs GPT-4 as a decision-making
component: it devises cognitive pathways for reasoning over human commonsense
and translates LLM decisions into low-level control. It demonstrates that LLMs
\emph{can} make driving decisions, but relies on cloud inference and does not
report timing, so it cannot close a real-time loop. DriveGPT4
\cite{xu2023drivegpt4} produces interpretable end-to-end outputs via a
vision-language model, prioritizing explainability over latency. A broader line
of work applies LLMs to scene understanding, visual question answering, or
offline/single-step planning; none addresses the latency required for
high-speed, safety-critical control. Our positioning: we treat the LLM as a
low-frequency tactical layer and show that a 1.5B local model is fast enough
(0.59~s) for real-time use --- an order of magnitude faster than the cloud-API
paradigm --- while safety-critical tracking remains rule-based.

\subsection{Hierarchical Control and Safety}
Classical driving stacks separate motion planning from trajectory tracking
(pure pursuit, Stanley, LQR, MPC) and enforce safety with control barrier
functions (CBF), QP-based filters, or set-based reachability. We adopt a
pure-pursuit tracking controller and a lightweight geometric \emph{signed
shield} that brakes when the ego is alongside the lead with insufficient
lateral clearance. We keep the shield rule-based and report, as a negative
result, that a full CBF-QP formulation underperformed it in a 2~m stress test
(Sec.~\ref{sec:neg}).

\subsection{Edge LLMs and Efficient Fine-Tuning}
Efficient deployment uses quantization and parameter-efficient fine-tuning such
as LoRA \cite{hu2021lora}. We fine-tune Qwen2.5-1.5B \cite{qwen2025qwen25} with
bf16 LoRA on a consumer 8~GB GPU, providing direct evidence that a tactical LLM
can run at the edge without cloud dependence.

\subsection{Latency in Real-Time Systems}
Networked control and robot action chunking address decision latency through
predictive compensation or coarser action granularity. We transfer action
chunking to high-speed driving and, unlike prior work, \emph{quantify} the
tolerable staleness (0.5~s) and the failure point (0.8~s), motivating our
latency reduction from 1.08~s to 0.59~s.

\subsection{Simulation and Model-Based Data}
We evaluate in CARLA \cite{dosovitskiy2017carla} on a closed Town04 ring. Our
counterfactual labeling uses the simulator as a ground-truth dynamics model for
a shallow action search, connecting to model-based reinforcement learning and
expert iteration; a learned neural world model to replace this simulator is
future work.

\section{Method}
\subsection{Problem Formulation}
We study a two-vehicle overtaking task on a closed ring. At each step the ego
observes a state consisting of ego and lead speeds, signed and absolute
longitudinal gap, distance to the corner apex, track curvature,
time-to-collision, current lateral offset and lateral separation, the
racing-line offset, and the corner's inside side (left/right). The tactical
decision is a structured intent
$a = (\text{action}, \text{extra\_inside\_m}, \text{speed\_factor})$ with
$\text{action}\in\{\text{overtake},\text{hold}\}$,
$\text{extra\_inside\_m}\in[0,3.5]$~m (additional lateral offset toward the
corner inside), and $\text{speed\_factor}\in[0.9,1.1]$. The objective is to
complete a pass while respecting a minimum lateral clearance and avoiding
collision.

\subsection{Layered Architecture}
The system decomposes into two layers with distinct timescales
(Fig.~\ref{fig:arch}).

\begin{figure}[t]
\centering
\includegraphics[width=\columnwidth]{architecture.pdf}
\caption{Two-timescale layered architecture. A low-frequency LLM tactical layer
issues overtake/hold intents; a 20~Hz rule-based control layer executes them
under a geometric safety shield. Decision latency is absorbed into the 1~Hz
cadence, not the 20~Hz control loop.}
\label{fig:arch}
\end{figure}

\textbf{Tactical layer (LLM, $\sim$1~Hz).} The fine-tuned LLM maps the
observation (plus the previous decision, which disambiguates ``start an
attack'' from ``already attacking'') to the intent $a$. The \emph{short-output}
variant omits the free-text \texttt{reason} field, cutting generated tokens by
roughly 60\% and latency by $\sim$45\% at a small accuracy cost.

\textbf{Control layer (20~Hz).} A pure-pursuit controller tracks a racing line
offset toward the corner inside. During an overtake it adds the attack offset
\texttt{extra\_inside\_m}, scaled by a spatial attack window centered on the
apex, and applies the speed factor. A \emph{signed shield} brakes when the
longitudinal gap is below 8~m and the lateral separation is below 2.2~m while
the ego is alongside, and caps the speed advantage over the lead at 6~m/s.

\subsection{Two Control-Layer Corrections}
\textbf{Attack-geometry correction.} The initial attack reference summed the
racing-line inside offset (3.3~m) and the attack offset (3.5~m) into a 6.8~m
lateral command that exceeded the available road width, driving the ego off the
road under variable and defensive leads. We reduced the racing-line inside
offset to 1.0~m, clamped the total offset to 4.5~m, and added a 0.5~m bias,
keeping the attack trajectory on the road without retraining the model.

\textbf{Inside-line hold.} The attack window ramps down after the apex and the
racing line returns to the outside; on tight corners the ego has not yet cleared
the lead at that point, so the retraction steered it back onto the lead's line
and caused side-swipe collisions. We pin the attack weight to unity while the
ego is still alongside the lead, holding the inside line until the pass
completes.

\subsection{Model and Counterfactual Training Data}
We fine-tune Qwen2.5-1.5B-Instruct with LoRA ($r{=}16$, $\alpha{=}32$, dropout
0.05) for two epochs with a completion-only loss (prompt tokens masked) on a
consumer 8~GB GPU. The training signal comes from two sources. First,
demonstrations from a rule-based expert driving in CARLA. Second, and more
importantly, \emph{counterfactual labels}: at sampled states we snapshot the
world and roll out five candidate actions (hold; inside 3.0~m; inside 3.5~m;
inside 3.5~m with late braking; outside 3.5~m) against the CARLA simulator for a
short horizon, and label the state with the lowest-cost action under a cost that
penalizes collision (weight 1000), sub-2.2~m clearance, and failed passes, while
rewarding longitudinal progress. This lets the model learn the \emph{optimal}
action in each state rather than the expert's demonstrated one, partially
escaping the behavior-cloning ceiling.

\subsection{Evaluation Protocol}
We build a closed Town04 ring with six corners (radii 45--208~m, in both left
and right directions) and define ten scenarios pairing five passable corners
with constant, variable, and defensive lead behaviors. Metrics are pass rate
(the ego reaches 2~m ahead during the attack window), collision rate, minimum
lateral separation (\texttt{minSep}), and decision latency. Pass rate is
reported with Wilson 95\% confidence intervals; \texttt{minSep} with bootstrap
confidence intervals.

\section{Experiments}
\subsection{Setup}
All experiments run in CARLA Town04 with a synchronous 20~Hz physics loop and a
ClearNoon setting. The LLM is served locally and queried at 1~Hz; the control
layer runs every tick. Latency is measured as end-to-end inference (prompt
encoding plus greedy decoding) on the local GPU.

\subsection{Latency--Accuracy Pareto}
\begin{table}[t]
\centering
\caption{Decision accuracy and end-to-end latency across model scales. The
1.5B-short variant is the Pareto-optimal point.}
\label{tab:pareto}
\begin{tabular}{lccc}
\toprule
Model & Accuracy & Latency & vs.\ 1.5B-short \\
\midrule
Qwen2.5-0.5B & 75.7\% & 0.86~s & 1.46$\times$ slower, $-10.4$~pt \\
Qwen2.5-1.5B & 87.8\% & 1.08~s & 1.83$\times$ slower, $+1.7$~pt \\
\textbf{Qwen2.5-1.5B-short (ours)} & \textbf{86.1\%} & \textbf{0.59~s} & --- \\
Qwen2.5-3B (4-bit) & 71.3\% & 2.02~s & 3.42$\times$ slower, $-14.8$~pt \\
Qwen2.5-7B (4-bit, zero-shot) & 69.6\% & 1.41~s & 2.39$\times$ slower, $-16.5$~pt \\
\bottomrule
\end{tabular}
\end{table}

\begin{figure}[t]
\centering
\includegraphics[width=\columnwidth]{latency_superiority.pdf}
\caption{(a) Decision latency and (b) accuracy--latency Pareto. The 1.5B
short-output model is 3.4$\times$ faster than 3B and 2.4$\times$ faster than 7B
while \emph{more} accurate.}
\label{fig:pareto}
\end{figure}

Two observations follow (Fig.~\ref{fig:pareto}, Table~\ref{tab:pareto}). First,
\emph{bigger is not better} here: the 3B and 7B models must be quantized to
4-bit to fit the 8~GB budget, which degrades their accuracy below the 1.5B
model, while their larger parameter count inflates latency. Second, the
\emph{short-output} ablation (dropping the \texttt{reason} field) cuts latency
from 1.08~s to 0.59~s ($-45\%$) at a cost of only 1.7 accuracy points, because
autoregressive decoding cost scales with the number of generated tokens. The
1.5B-short model is therefore the Pareto-optimal point: the fastest and nearly
the most accurate.

\subsection{Closed-Loop Overtaking (20 seeds, corner R=188)}
\begin{table}[t]
\centering
\caption{Closed-loop overtaking over 20 seeds on the R=188 corner. Median
clearance meets the 2.2~m target in all lead behaviors.}
\label{tab:closed}
\begin{tabular}{lcccc}
\toprule
Lead & Pass rate [95\% CI] & Collision & minSep med.\ & minSep min \\
\midrule
constant & 100\% [84,100] & 0\% & 2.87~m & 2.54~m \\
variable & 95\% [76,99] & 0\% & 3.13~m & 1.14~m \\
defensive & 95\% [76,99] & 5\% & 2.27~m & 1.04~m \\
\bottomrule
\end{tabular}
\end{table}

The median clearance meets the 2.2~m safety target in all three lead behaviors;
the minimum drops below 2~m only in rare tail cases (defensive and variable
leads). The defensive lead --- which moves toward the ego's chosen side --- is
the hardest, yielding one collision in twenty runs.

\subsection{Zero-Shot Generalization (10 scenarios, 10 seeds)}
Trained on only two corner geometries (one right, one left), the model
generalizes zero-shot across ten corner--lead scenarios (Table~\ref{tab:zero}).
Aggregate: \textbf{96/100 passes (96\%), 2 collisions (2\%)}. Seven scenarios
reach 100\% pass with zero collision; two defensive-lead scenarios reach 90\%;
the single tight-corner scenario (R=83) drops to 80\% with two collisions.

\begin{table}[t]
\centering
\caption{Zero-shot generalization across ten corner--lead scenarios (10 seeds
each).}
\label{tab:zero}
\begin{tabular}{llcccc}
\toprule
Corner & Lead & Pass & Collision & minSep med.\ & minSep min \\
\midrule
R=208 & constant & 100\% & 0 & 2.77~m & 2.71~m \\
R=208 & variable & 100\% & 0 & 3.30~m & 3.22~m \\
R=208 & defensive & 90\% & 0 & 2.33~m & 1.01~m \\
R=151 & constant & 100\% & 0 & 3.28~m & 2.83~m \\
R=151 & variable & 100\% & 0 & 2.79~m & 1.60~m \\
R=151 & defensive & 100\% & 0 & 2.37~m & 2.34~m \\
R=188 & constant & 100\% & 0 & 2.84~m & 2.01~m \\
R=188 & defensive & 90\% & 0 & 2.28~m & 0.65~m \\
R=83 & constant & 80\% & 2 & 2.74~m & 1.63~m \\
R=138 & variable & 100\% & 0 & 3.26~m & 3.24~m \\
\bottomrule
\end{tabular}
\end{table}

Two findings stand out. First, generalization is governed by \emph{corner
difficulty} (curvature) rather than direction: gentle corners pass in both
directions, while the tight right corner R=83 is marginal. Second, the defensive
lead is the hardest \emph{behavior}, consistently lowering clearance across
corners.

\subsection{Stale-Latency Boundary}
\label{sec:stale}
To justify the latency reduction, we measure closed-loop safety as a function of
injected decision delay. With 0.5~s of added staleness the pass rate and
clearance are essentially unchanged, confirming that a 0.59~s decision is within
the tolerable regime. At 0.8~s the shield begins to intervene more often and
safety degrades, marking the boundary beyond which the decision is stale. This
quantifies the enabling condition for the cloud-to-edge substitution: the
1.5B-short model sits \emph{inside} the tolerable regime, whereas a 2~s (3B) or
multi-second (cloud) decision would sit far outside it.

\subsection{Ablations}
\textbf{Short-output.} Dropping the \texttt{reason} field reduces latency from
1.08~s to 0.59~s ($-45\%$) at a 1.7-point accuracy cost (Table~\ref{tab:pareto}),
confirming that the free-text explanation is the dominant decoding cost.

\textbf{Control corrections.} Without the attack-geometry correction, the ego
drives off the road under variable and defensive leads (0\% pass, collisions);
with it, those scenarios recover to 95--100\% pass and zero collision. Without
the inside-line hold, the tight left corner R=138 collides in 6/6 runs
(minSep $\approx$1.9~m); with it, clearance rises above 3~m and collisions
disappear. Both corrections are therefore necessary, not incidental.

\textbf{Counterfactual labels.} Replacing counterfactual labels with pure
behavior cloning from the rule expert makes multi-scenario fine-tuning regress
(see Sec.~\ref{sec:neg}), indicating that the labels, though few, carry
disproportionate signal.

\subsection{Negative Results}
\label{sec:neg}
\textbf{Multi-scenario fine-tuning regresses.} Fine-tuning on a ten-scenario
dataset collected from the same rule expert \emph{reduced} closed-loop pass rate
(left-turn scenarios fell from 100\% to 0--67\%); merging the new data with the
original two-corner data did not recover the zero-shot baseline. Root cause:
the rule expert only \emph{holds} on hard scenarios, so its 607 demonstrated
actions (mostly ``hold'') overwhelm the 47 counterfactual labels, teaching
conservatism. This is a concrete demonstration that behavior cloning from a
weak expert sets a ceiling, and that the counterfactual-labeling ratio must be
much higher for multi-scenario scaling.

\textbf{CBF-QP vs.\ geometric shield.} A full CBF-QP safety filter
underperformed a plain if-else shield in a 2~m stress test, consistent with
prior findings that hand-tuned geometric gates can be more robust than
optimization-based filters for narrow, well-understood tasks. We therefore
retain the lightweight geometric shield.

\section{Discussion}
\textbf{Limitations.} Our study is confined to a single map (Town04), a single
task (overtaking), and simulation-only evaluation. We have not validated
transfer to other maps, weather conditions, or interactive multi-vehicle
traffic, and we have no real-world (closed-course or on-road) deployment; the
sim-to-real gap --- sensor noise, latency jitter, and unmodeled dynamics ---
remains open. The tactical action space is intentionally coarse (three
scalars), sufficient for overtaking but not for general driving, so the approach
is a building block rather than a complete driving system. Our advantage is
empirical and system-level rather than algorithmic: the base model and
fine-tuning method are off the shelf, and ``smaller is faster'' is expected ---
the contribution is the \emph{quantification} of the operating point and the
closed-loop evidence that it suffices. Finally, the behavior-cloning ceiling
means the model inherits the weaknesses of its data source on hard scenarios
(tight corners, defensive leads), which remain its weakest cases.

\textbf{Implications.} The central message is that the latency question, not
model capability, is the binding constraint for embedding LLMs in real-time
driving, and that a two-timescale architecture resolves it. The
counterfactual-labeling mechanism also suggests a path beyond the
behavior-cloning ceiling, connecting to model-based reinforcement learning and
expert iteration.

\textbf{Future work.} (1) Train a lightweight neural \emph{world model} to
replace the CARLA simulator in counterfactual search, potentially reducing
decision latency toward the millisecond scale; (2) scale up counterfactual
labeling to escape the weak-expert ceiling and enable multi-scenario
fine-tuning; (3) extend to multiple maps, weather, and interactive traffic, and
toward real-vehicle validation.

\section{Conclusion}
We showed that a locally fine-tuned 1.5B LLM can serve as a low-frequency
tactical decision layer for high-speed overtaking at 0.59~s latency --- an order
of magnitude faster than the cloud-API paradigm --- while a 20~Hz rule-based
controller guarantees safety. The system achieves 95--100\% pass rate with
median clearance $\geq$2.2~m and generalizes zero-shot across ten scenarios. We
quantified the stale-latency boundary that justifies the reduction, and reported
honest negative results that bound the approach and point to future work.

\section*{Acknowledgment}
The authors thank the anonymous reviewers for their constructive feedback.
(Optional; fill before submission.)

\begin{thebibliography}{5}
\bibitem{sha2023languagempc}
H.~Sha, Y.~Mu, Y.~Jiang, L.~Chen, C.~Xu, P.~Luo, S.~E.~Li, M.~Tomizuka,
W.~Zhan, and M.~Ding, ``LanguageMPC: Large language models as decision makers
for autonomous driving,'' \emph{arXiv:2310.03026}, 2023.

\bibitem{xu2023drivegpt4}
Z.~Xu, Y.~Zhang, E.~Xie, Z.~Zhao, Y.~Guo, K.-Y.~K.~Wong, Z.~Li, and H.~Zhao,
``DriveGPT4: Interpretable end-to-end autonomous driving via large language
model,'' \emph{arXiv:2310.01412}, 2023.

\bibitem{hu2021lora}
E.~J.~Hu, Y.~Shen, P.~Wallis, Z.~Allen-Zhu, Y.~Li, S.~Wang, L.~Wang, and
W.~Chen, ``LoRA: Low-rank adaptation of large language models,'' in
\emph{Int. Conf. Learning Representations (ICLR)}, 2022.

\bibitem{dosovitskiy2017carla}
A.~Dosovitskiy, G.~Ros, F.~Codevilla, A.~Lopez, and V.~Koltun, ``CARLA: An open
urban driving simulator,'' in \emph{Conf. Robot Learning (CoRL)}, 2017.

\bibitem{qwen2025qwen25}
Qwen Team, ``Qwen2.5 technical report,'' \emph{arXiv:2412.15115}, 2024.
\end{thebibliography}

\end{document}
